\begin{lemma}[Normalizing Situation 7.2 Case 3]
Fix a minimum hierarchy counterexample and a raw
\noderef{HierarchySituationTwoCaseThreeData} \(D\).  Thus
\(p(T)=\ell\ge2\), and its color \(\alpha\) is missing at the terminal
vertex \(y_p\) and at an explicitly displayed vertex
\(v\prec y_{p-1}\).  Then there is a
\noderef{HierarchyReplacementPresentation} \(P\) on the same exact
terminal tree \(T\) for which Case 3 has a
\noderef{HierarchySituationTwoCaseThreeNormalizedData}.  Equivalently,
after passing to an actual corresponding same-rung maximum ETT and
transporting the complete minimum context when necessary, condition (20)
may be assumed.  The conclusion is independent of whether the last-edge
color \(\theta\) is missing at \(y_{p-1}\).
\end{lemma}
\begin{proof}
Write \(H=C.\mathit{hierarchy}\), let \(q\) denote its divider count and
also use \(q\) for the last finite divider index, and put
\[
 A=T_{n,q}^{*}=\noderef{HierarchyStarredStage}(H,q),\qquad
 \ell=\noderef{HierarchyFullTailLength}(C).
\]
For \(0\le i\le\ell\), write
\(T(y_i)=\noderef{HierarchyFullTailPrefix}(C,i)\).  In particular set
\[
 B=T(y_{p-2})=T(y_{\ell-2}),
 \qquad B^+=T(y_{p-1})=T(y_{\ell-1}).                           \tag{1}
\]
The common-tail data give \(\ell\ge2\).  Hence all indices in (1) are
defined, \(A\preccurlyeq B\prec B^+\prec T\), and \(T-B\) consists of
precisely the last two individual edges.  The position equations in
\noderef{HierarchySituationTwoCommonTailData} identify those last two
positions with the displayed penultimate and terminal additions; no count
of endpoint pairs is being substituted for a count of edge identities.

First derive the prefix-elementarity premise needed for counting.  Apply
\noderef{HierarchyCounterexamplePenultimateElementary} using properness,
exact one-uncoloredness, the six lower-rung assertions, fewer-divider
induction, lexicographic minimality, and the \(2n+11\) estimate, all of
which are fields of
\noderef{HierarchyMinimumCounterexampleContext}.  We obtain
\[
                       V(B^+)\text{ is elementary}.            \tag{2}
\]

We now apply \noderef{NineN} at the exact index \(i=\ell-2<\ell\).
Use the notation of that lemma:
\[
 X=(\overline\varphi_n(B)\setminus M_0)\setminus U_{\ell-2},
 \quad D=D_{n,q},\quad
 \Gamma^q=\bigcup_{\eta\in D}\Gamma^q_\eta,                  \tag{3}
\]
where \(M_0\) is the set of colors missing on
\(T_{n,0}^{*}-V(T_n)\), and \(U_{\ell-2}\) is the set of colors on the
exact prefix-tail edges from \(A\) through \(B\).  The hypotheses of
\noderef{NineN} are (2) and the large-set field of the context.  Its first
conclusion gives
\[
                         |X|\ge2n+11.                           \tag{4}
\]

We split according to the size of
\(X\setminus(\Gamma^q\cup D)\).

Suppose first that
\[
                 |X\setminus(\Gamma^q\cup D)|\ge5.             \tag{5}
\]
Passing from the edge interval ending at \(B\) to the full tail adds
exactly the two final edge identities.  Their two colors delete at most
two members of the set in (5).  Consequently at least three colors remain.
Choose one and call it \(\beta\).  Directly from (3) and the choice,
\[
 \beta\in\overline\varphi_n(B),
 \quad \beta\notin M_0,
 \quad \beta\notin\varphi_n\langle T-A\rangle,
 \quad \beta\notin D_{n,q}\cup\Gamma^q.                       \tag{6}
\]
When \(q>0\), (6) is the first alternative required in (21).

At \(q=0\), the required statement is stronger: it excludes \(\beta\)
from the raw \(D_n\), not only from \(D_{n,0}\).  We verify it rather than
silently identifying the two sets.  By
\noderef{HierarchyStoredColorsAt},
\[
                  D_{n,0}=D_n\setminus\overline\varphi_n(A).   \tag{7}
\]
If \(\beta\in D_n\), then \(\beta\notin D_{n,0}\) from (6) and (7)
would force \(\beta\in\overline\varphi_n(A)\).  A raw stored color is
not missing on \(T_n\), by the defining subtraction in
\noderef{TashkinovStoredColors}.  Since at divider zero \(A\) is the
hierarchy core, every remaining vertex of \(A\) lies in
\(T_{n,0}^{*}-V(T_n)\).  Thus \(\beta\in M_0), contradicting (6).
Therefore
\[
                 \beta\notin D_n\cup\Gamma^0                 \tag{8}
\]
when \(q=0\).  This is precisely the raw-set upgrade at paper line 4210.

Now suppose instead that
\[
                 |X\setminus(\Gamma^q\cup D)|\le4.             \tag{9}
\]
The second conclusion of \noderef{NineN} supplies an injectively indexed
family \((\eta_h)_{h\in\operatorname{Fin}7}\) such that, for every \(h\),
\[
 \eta_h\in D\cap\overline\varphi_n(B),
 \qquad
 (\Gamma^q_{\eta_h}\cup\{\eta_h\})\cap U_{\ell-2}=\varnothing.\tag{10}
\]
These seven three-color blocks are pairwise disjoint.  Indeed, injectivity
makes their stored colors distinct; the `disjointPairs' field of
\noderef{GoodHierarchy} separates their Gamma pairs; and
\noderef{GoodHierarchyReservedDisjoint} says that the full Gamma union is
disjoint from \(D\).  Each of the two final edge colors can therefore meet
at most one block.  Since there are seven blocks, choose \(h\) whose entire
block meets neither final edge color, and put \(\beta=\eta_h\).  Combining
this choice with (10) gives
\[
 \beta\in D_{n,q}\cap\overline\varphi_n(B),
 \qquad
 (\Gamma^q_\beta\cup\{\beta\})
       \cap\varphi_n\langle T-A\rangle=\varnothing.             \tag{11}
\]

Thus in both size branches we have selected a color \(\beta\) satisfying
\[
 \beta\in\overline\varphi_n(B)
       \setminus\varphi_n\langle T-A\rangle,                   \tag{12}
\]
together with exactly one of: the positive-divider exclusion in (6), the
divider-zero raw exclusion (8), or the stored-color alternative (11).
This is condition (21), with every exact edge of the full tail accounted
for.

We next turn (12) into a normalized Case 3 presentation.  If
\(\beta\in\overline\varphi_n(y_p)\), take the identity replacement
presentation.  Concretely, reuse the original input, series, ETT, base,
hierarchy, good-hierarchy certificate, counterexample, and all fields of
the minimum context.  The identity equalities fill every field of
\noderef{CorrespondingExtendedTashkinovTree}, so this is an actual
\noderef{HierarchyReplacementPresentation}, not a bare assertion that no
recoloring occurred.  Choose from (12) a vertex of \(B\) missing
\(\beta\).  Its displayed index is strictly before \(y_{p-1}\) by (1),
while \(\beta\) is also missing at \(y_p\).  Equations (6), (8), or (11),
together with (12), fill the inherited raw Case 3 fields and all fields of
\noderef{HierarchySituationTwoCaseThreeNormalizedData}.  This proves the
theorem in this branch.

Assume henceforth that \(\beta\notin\overline\varphi_n(y_p)\).  The raw
Case 3 hypothesis gives
\(\alpha\in\overline\varphi_n(y_p)\), so \(\alpha\ne\beta\).  Let
\[
 Q=P_{y_p}(\alpha,\beta,\varphi_n),
 \qquad \sigma_n=\varphi_n/Q.                                 \tag{13}
\]
Properness from the context and
\noderef{TwoColorComponentSpanningPath} allow us to traverse \(Q\) as a
finite simple alternating path starting at \(y_p\).

If \(Q\) has no vertex in \(B^+\), apply
\noderef{HierarchySituationTwoOutsidePenultimateReconstruction}.  Its
candidate hypotheses are exactly (12) and the appropriate one of (6),
(8), and (11), and its terminal-present hypothesis is the present branch
assumption.  The cited lemma returns an actual replacement presentation
on \(T\) and normalized datum with color \(\beta\).  Notice that this use
treats \(Q\) as outside \(B^+\) only; since \(y_p\in Q\cap V(T)\), it does
not invoke \noderef{HierarchyLEMStable} on the full \(T\).

It remains to suppose that \(Q\) meets \(B^+\).  Let \(u\) be the first
vertex of \(B^+\) encountered while traversing the path from \(y_p\).  If
\(u=y_{p-1}\), form the exact tree-sequence
\[
                      T'=B^+\cup Q[u,y_p];                     \tag{14a}
\]
otherwise form
\[
                      T'=B\cup Q[u,y_p].                       \tag{14b}
\]
In either formula the union means: retain the indicated ordered prefix
entry by entry, then append in traversal order the individual path-edge
identities and new path vertices from \(u\) to \(y_p\).  Firstness of \(u\)
and simplicity of the spanning path show that each appended vertex is new;
looplessness and the edge/vertex incidence equations show that every append
is an exact tree-sequence extension.  Both \(\alpha\) and \(\beta\) are
missing on \(B\): the former because its raw witness precedes \(y_{p-1}\),
and the latter by (12).  The alternating path can therefore be appended by
\noderef{TashkinovAugmentation}s under \(\varphi_n\).  Here is the exact
route argument.  In the PE terminal constructor of
\noderef{TashkinovSeriesStep}, \noderef{TashkinovClosure} gives a route
from \(T_n\) to the final closure.  In the SE or RE constructor it records
the individual connecting edge \(f_n\), its one-edge seed, and the route
from that seed.  Since \(B\), \(B^+\), and \(T\) are exact prefixes of the
same final closure, \noderef{TashkinovAugmentationSequenceRestrict}
retains the appropriate route up to the base used in (14).  Append the
edges of \(Q[u,y_p]\) in their traversal order.  For the first append its
edge color is one of \(\alpha,\beta\), both of which are tree-missing on
the starting prefix.  Inductively, alternation shows that the next edge
color is still missing either at its original witness in \(B\), or at the
new endpoint of the preceding path edge.  Firstness of \(u\) ensures that
this endpoint is new.  This proves a Tashkinov-augmentation sequence to
the exact tree \(T'\).

Apply \noderef{Td} to the original ETT with \(\sigma=\varphi_n\):
properness and exact one-uncoloredness are fields of the minimum context,
stability is reflexive, and the lower-rung package is the context's six
assertions.  Thus \(\noderef{ModColoring}(\mathcal E,\varphi_n)\).
The route just constructed satisfies the PE or exact SE/RE eligibility
branch in its universal clause, and therefore gives an actual
\noderef{CorrespondingExtendedTashkinovTree} on \(T'\).
\noderef{MaximumPropertyTransport} transfers the original maximum
property to it.  This accounts separately for the ETT and maximum data
rather than treating them as automatic consequences of adjoining a path.

The corresponding-ETT certificate identifies all ladder trees, extension
kinds, connecting edges, and stored data.  Therefore the core, closure,
support position, PE clause, and non-PE clause of
\noderef{HierarchyBase} copy to a base on the new ETT.  Define a
\noderef{TashkinovHierarchy} with the old divider count and all old
nonterminal stages, and with terminal stage \(T'\).  Its nesting and
strictness follow from the exact-prefix construction (14); its divider
positions and divider-stage equalities are unchanged.  Thus all computed
stored sets and all stage-missing sets are unchanged.  We retain the old
Gamma pairs.  Their cardinality, pairwise disjointness, unique-refresh,
PE-base, and minus-closed fields copy literally.  In the two branches
below we check the sole remaining issue, namely avoidance from the new
terminal segment-color sets.

By the choice of the first intersection, the new alternating segment
begins no later than the final path tail of \(T\), and the definition in
(14) discards at least one old path-tail step.  Hence
\[
                         p(T')<p(T).                            \tag{15}
\]

The candidate \(\beta\) is outside \(\Gamma^q\) in every alternative.
This is explicit in (6) and (8).  In (11), it follows from
\(\beta\in D_{n,q}\) and the disjointness conclusion of
\noderef{GoodHierarchyReservedDisjoint}.  We now consider the three
possible positions of the original \(\alpha\).

If \(\alpha\notin\Gamma^q\), copy all old hierarchy stages and Gamma
pairs to \(T'\).  Their avoidance clauses survive because both colors used
to append (14) are outside \(\Gamma^q\).  Together with the copied fields
checked above, this gives a good hierarchy with the same divider count.
The two old \(\alpha\)-missing vertices, the
earlier one in \(B\) and \(y_p\) on the appended path, both lie in \(T'\),
so \(T'\) is not elementary.  We have therefore constructed a genuine
hierarchy counterexample \(C'\).

Suppose next that
\(\alpha\in\Gamma^q\cap\varphi_n\langle T-A\rangle\).  The last edge is
not colored \(\alpha\), since it is incident with the vertex \(y_p\) at
which \(\alpha\) is missing.  Hence an \(\alpha\)-colored tail edge occurs
already before \(B^+\).  Write \(\alpha\in\Gamma^q_\eta\).  The avoidance
field of \noderef{GoodHierarchy}, evaluated at the first such segment,
puts the stored color \(\eta\) in
\(\overline\varphi_n(B)\).  Together with
\(\beta\in\overline\varphi_n(B)\) and
\(\beta\notin\Gamma^q\), this verifies the same avoidance clauses for the
hierarchy on \(T'\): every newly appended \(\alpha\)-edge occurs after the
first \(\eta\)-missing position, and the other appended color is not a
Gamma color.  The preceding copied-field argument now constructs the
complete good hierarchy.  Again the old \(\alpha\)-missing vertices make its
terminal set non-elementary, so this branch also produces a genuine
counterexample \(C'\) with the same rung and divider count.

In each of these first two subbranches, (15) strictly decreases the first
coordinate of
\noderef{HierarchyCounterexampleMeasure}.  Therefore
\[
 \noderef{HierarchyCounterexampleLess}
   (\noderef{HierarchyCounterexampleMeasure}(C'),
    \noderef{HierarchyCounterexampleMeasure}(C))
\]
holds without inspecting the second coordinate.  This contradicts the
lexicographic-minimality field of the supplied minimum context.  Thus both
subbranches are eliminated.

The only remaining case is
\[
          \alpha\in\Gamma^q
             \setminus\varphi_n\langle T-A\rangle.             \tag{16}
\]
Choose \(\eta\in D_{n,q}\) with
\(\alpha\in\Gamma^q_\eta\).  The avoidance field of
\noderef{GoodHierarchy} puts \(\alpha\) in the corresponding
\noderef{HierarchyStageMissingColors}: on the last ordinary stage when
\(q>0\), and on \(T_n\) when \(q=0\).  In particular \(\alpha\) is missing
on \(A\).  Both \(\alpha\) and \(\beta\) are missing on \(B^+\).

We verify the conditional order premise of \noderef{ChangeLemma}.  If
\(\beta\) is also missing on \(A\), the premise's antecedent is false,
because both colors are missing there.  Otherwise the first
\(\alpha\)-missing vertex lies in \(A\), while the first
\(\beta\)-missing vertex lies after \(A\) and no later than its witness in
\(B\).  Thus the former precedes the latter.  Equation (16) says that
\(\alpha\) is unused on every exact edge after \(A\), so in particular it
is unused up to that first \(\beta\)-missing vertex.  This is precisely the
required order-and-unused implication.  The lower-rung, fewer-divider,
same-rung minimality, prefix-elementarity, and large-set inputs are the
corresponding fields of the minimum context and (2).  Finally, the
stage-missing conclusion above supplies the positive-divider case of the
lemma, or its \(T_n\)-missing divider-zero case.  Consequently
\noderef{ChangeLemma} gives
\[
 P_{v_\alpha}(\alpha,\beta,\varphi_n)
   =P_{v_\beta}(\alpha,\beta,\varphi_n),                     \tag{17}
\]
where \(v_\alpha\) and \(v_\beta\) are the respective first missing
vertices.  Both belong to \(B^+\).

The component in (17) is not \(Q\).  If
\(v_\alpha=v_\beta\), that vertex misses both colors and hence is an
isolated vertex of the two-color incidence graph, whereas \(Q\) also
contains the distinct terminal vertex \(y_p\).  If they are distinct,
the common component in (17) contains one vertex missing \(\alpha\) and
one missing \(\beta\); equality with \(Q\) would add the third distinct
degree-at-most-one vertex \(y_p\), which misses \(\alpha\).  Properness,
\noderef{TwoColorIncidenceDegreeAtMostTwo}, and
\noderef{TwoColorComponentSpanningPath} make every such component a
finite path, which has only two endpoints.  Both cases are impossible.
The dichotomy in \noderef{TwoColorComponentsDisjoint} therefore makes
\(Q\) disjoint from the component in (17).

Now let \(x\in V(B^+)\).  Elementarity (2) says that if \(x\) misses
\(\alpha\), then \(x=v_\alpha\), and if it misses \(\beta\), then
\(x=v_\beta\).  Hence a vertex of \(Q\cap V(B^+)\) misses neither color.
Properness gives it one incident edge of each color, both in \(Q\), and
the swap exchanges those two edges together; its missing set is unchanged.
At a vertex outside \(Q\), every incident edge is unchanged.  Thus every
vertex of \(B^+\) has the same missing set under \(\sigma_n\) as under
\(\varphi_n\).

Equations (12) and (16) say that neither \(\alpha\) nor \(\beta\) colors
an exact full-tail edge.  Hence the swap fixes every individual displayed
edge of \(T\), including the two final edges.  The first missing position
of every stored color is also fixed: a stored \(\alpha\) or \(\beta\)
already has its first missing witness in \(B\), while all other colors are
unchanged.  These facts verify, respectively, the full-tail edge,
penultimate missing-set, and stored-first-position hypotheses of
\noderef{StableChangeCasesOneTwoFullTail}.  Its Cases 1/2 premise is
supplied with \(\alpha\) in the stage-missing slot, as just obtained from
(16); no such premise is imposed on the candidate \(\beta\).  The cited
theorem returns a weakly stable mod-coloring, an actual corresponding
same-rung maximum ETT on the exact \(T\), and a rebuilt good hierarchy with
the same stages, stored sets, and Gamma pairs.

Under the swap (13), \(\beta\) becomes missing at \(y_p\), while its
missing witness in \(B\) is preserved by the \(B^+\) missing-set equality.
Thus the returned terminal tree is non-elementary and gives a genuine
replacement counterexample.  Build the replacement minimum context as in
\noderef{HierarchySituationTwoOutsidePenultimateReconstruction}: retain
the universal six lower-rung and fewer-divider fields, transport
lexicographic minimality across equality of the two measure coordinates,
and transport the \(2n+11\) estimate across ladder-tree stability.  This
fills a \noderef{HierarchyReplacementPresentation}.  Finally (12) and
the appropriate alternative among (6), (8), and (11), transported across
the full-tail edge equality and copied Gamma data, fill the normalized
datum with color \(\beta\).

Every branch that survives the shorter-\(T'\) contradiction has now
returned both an actual replacement presentation and a complete normalized
Case 3 datum.  The argument never uses any assertion about whether
\(\theta\) is missing at \(y_{p-1}\), so the normalization is common to
Subcases 3.1 and 3.2.
\end{proof}
